\uuid{4G7c}
\titre{ Réseau entièrement linéaire : équivalence avec un neurone simple }
\niveau{L3}
\module{Analyse de données}
\chapitre{Réseaux de neurones}
\sousChapitre{Fonctions d'activation, expressivité}
\theme{Activation linéaire, composition de fonctions linéaires, neurone unique}
\auteur{Maxime Nguyen}
\datecreate{2026-03-19}
\organisation{AMSCC}
\difficulte{3}

\contenu{
\texte{
  On considère un réseau $1 \to 3 \to 1$ dont toutes les activations sont linéaires (identité).
}
\begin{enumerate}

\item
  \question{Ce réseau est-il différent d'un neurone simple $1 \to 1$ ? Justifier. La réponse change-t-elle par rapport au cas $1 \to 1 \to 1$ avec $h = 1$ ?}
  \reponse{
    Non. Toute composition de fonctions linéaires est linéaire.
    Formellement, avec $W^{[1]} \in \mathbb{R}^{1 \times 3}$, $W^{[2]} \in \mathbb{R}^{3 \times 1}$ :
    \[
      \hat{y} = W^{[2]}\bigl(W^{[1]} x + b^{[1]}\bigr) + b^{[2]}
             = \underbrace{W^{[2]} W^{[1]}}_{\tilde{w} \in \mathbb{R}} x + \underbrace{W^{[2]} b^{[1]} + b^{[2]}}_{\tilde{b} \in \mathbb{R}}.
    \]
    Le réseau équivaut donc à un neurone simple $\hat{y} = \tilde{w} x + \tilde{b}$.
    La réponse est la même que pour $h = 1$ : quelle que soit la largeur de la couche cachée,
    un réseau entièrement linéaire n'a pas plus de capacité d'expression qu'un modèle linéaire simple.
  }

\end{enumerate}
}
